\subsection{More on Compactness}
\thmp{Extreme Value Theorem}{
    Let $X$ be a compact space. Let $f: X \to \R$ be a continuous function. Then $f$ achieves its minimum and maximum values. i.e. $\exists a,b \in X$ s.t. $f(a) \leq f(x) \leq f(b)$ for all $x \in X$
}{
    $f(X)$ is a compact subset of $\R$. It is closed and bounded. Let $s=\sup f(X)$

    Consider the family of intervals $\{(-\infty, t) \mid t < s\}$

    If $s \notin f(X)$

    This is an open covering of $f(X)$ which has no finite subcover, since 
}

\defn{Uniformly Continuous}{
    Let $(X,d_X)$, $(Y,d_Y)$ be metric spaces.
    
    A function $f: X \to Y$ is uniformly continuous if $\forall \epsilon > 0$, $\exists \delta > 0$ s.t. 
    \[d_X(x_1,x_2) < \delta \implies d_Y(f(x_1),f(x_2))<\epsilon\]
}

\thm{Uniform Continuity Theorem}{
    Let $X,Y$ be metric spaces, and suppose $X$ is compact.

    Then any continuous function $f: X \to Y$ is uniformly continuous.
}

\thmp{}{
    Let $\mathcal{C}$ be an open covering of a compact metric space $X$. Then there is some $\delta > 0$ s.t. every subset of $X$ with diameter $< \delta$ is contained in some set $C \in \mathcal{C}$

    $\delta$ is called a Lebesgue number for the cover $\mathcal{C}$
}{
    For each $x \in X$, let $C_x$ be a set in $\mathcal{C}$ containing $x$.
    There is a ball $B(x,r_x) \subset C_x$ for some $r_x > 0$.

    Consider the open covering of $X$ give by the set of balls
    \[\{B(x,\frac{r_x}{2}) \mid x \in X\}\]
    $X$ is compact, so take a finite subcover
    \[B(x_1,\frac{r_1}{2}),\dots,B(x_n,\frac{r_n}{2})\]
    Let $\delta=\min(\frac{r_1}{2},\dots,\frac{r_n}{2}) > 0$
    
    Let $A \subset X$ have diameter $< \delta$. Then for $z \in A$ is contained in a ball $B(x_i, \frac{r_i}{2})$.

    For any $y \in A$, $d(y,z) < \delta$, so $d(x_i,y)<d(x_i,z)+d(z,y)<\frac{r_i}{2}+\delta \leq r_i$.

    Thus $A \subseteq B(x_i,r_i) \subseteq C_{x_i}$
}

\begin{thmpf}[Uniform Continuity Theorem]:

    Let $\epsilon > 0$. Consider open covering of $Y$ given by
    \[\{B(y,\frac{\epsilon}{2}) \mid y \in Y\}\]
    The sets
    \[\{f^{-1}(B(y,\frac{\epsilon}{2})) \mid y \in Y\}\]
    give an open covering of $X$.

    Let $\delta$ be a Lebesgue number for this covering. 
    
    If $x_1,x_2 \in X$ s.t. $d(x_1,x_2) < \delta$, diam$\{x_1,x_2\} < \delta$, so $x_1,x_2$ both lie in some 
    
    $f^{-1}(B(y,\frac{\epsilon}{2})) \implies f(x_1),f(x_2) \in B(y, \frac{\epsilon}{2}) \implies d(f(x_1),f(x_2)) < \epsilon$.
\end{thmpf}

\subsubsection{The Cantor Set}
Let $C_0=[0,1]$

Let $C_1=C_0$ minus middle third of interval

Let $C_2=C_1$ minus middle third of each component interval of $C_1$

In general,

$C_i$ is a union of closed intervals

$C_{i+1}$ is the set obtained by deleting "middle third" interval of each connected component of $C_i$

The Cantor set is $C=\displaystyle\bigcap_{i=0}^\infty C_i$

No endpoint of an interval in any $C_i$ is deleted. In particular, $0,1$ all in $C$, $3^{-n} \in C$ always.

Each $C_i$ is closed in $\R$ finite union of closed intervals thus $C$ is closed, and bounded by $0,1$, so it is compact.

It contains many embedded copies of itself.

\defn{Isolated}{
    A point $x$ in a space $X$ is called isolated if $\{x\}$ is open in $X$.
}

Cantor set contains no isolated points.

\thm{}{
    Let $X$ be any compact Hausdorff space.

    If $X$ has no isolated points, then it is uncountable.
}

$C$ shows up all the time.

Let $\{X_j\}_{j \in \N}$ be a family of spaces s.t. each $X_j$ is finite, discrete, and non-empty.
\[\prod_{j \in \N}X_j\]
is always homeomorphic to Cantor set.

$\{0,1\}^\omega$ is homeomorphic to Cantor set.

Can specify a point in Cantor set by specifying which components of each $C_i$ it belongs to, same as a sequence of left right choices.

\thm{}{
    Let $X$ be any space satisfying the following:
    \begin{itemize}
        \item $X$ is compact
        \item $X$ is metrizable
        \item $X$ has no isolated points
        \item $X$ is totally disconnected
    \end{itemize}
    $X$ is homeomorphic to the Cantor set.
}

\thm{}{
    Let $X$ be any compact metrizable space. Then $X$ is homeomorphic to a quotint $C \diagup \sim$ where $C$ is the Cantor set.
}
